\section{Question and principal results}\label{sec:question}
The geometric motivation is a central starting or meeting region, opposed upward and downward routes, outward spreading and a curved return, possibly with transfer elsewhere, while the native supports remain stationary. A correspondence argument must connect that picture to the actual state variables and exchange laws. Storage gained by a receiver is an energetic statement; crossing a specified spatial boundary is a separate geometric and dynamical statement.

The inherited candidate already supplies substantial structure: positioned native support lines, a registered carrier curve with two finite host contacts, reciprocal native feedback, a reversible receiver and a matched work account \cite{AX,CC,RX,RR}. The present paper asks how far those structures realize the intended route. It combines two accepted internal investigations \cite{CR,SO} into one argument at the unchanged registration.

The first result is a temporal distinction. From empty carrier and receiver preparation with nonzero registered drive, positive receiving occurs before activity at the carrier point nearest the central origin. A related conclusion holds at diagnostic cuts with a positive distance from the complete source-support union. The regional account explains why this does not violate conservation: the receiver couples to the same distributed interfaces as the native drive, and complementary work must be retained whenever the diagnostic region omits part of those interfaces.

The second result is spatial. For every admitted static geometry parameter and both hands, every inherited circle lies at distance greater than $0.0594\wh$ from each central opening boundary and its individual convex hull, where $\wh$ is the common host width. This is an exact uniform lower bound, not a sampled intersection test. Only one of the three geometric circles is instantiated dynamically. The result excludes direct crossings by that unchanged carrier, while leaving room for an explicitly added spatial connection.

Sections~\ref{sec:model}--\ref{sec:launch} establish the model, radial account and receiving comparison. Section~\ref{sec:separation} proves the uniform separation. Section~\ref{sec:transport} combines the geometric result with direction and information limits. Throughout, $R_B$ denotes a diagnostic \emph{ball radius}, whereas $r_\perp(P)$ denotes a \emph{cylindrical radius}. Appendix~\ref{sec:crosswalk} records this and the other editorial symbol changes explicitly.
